\begin{frame}[allowframebreaks]{Summary}
  \begin{itemize}
    \item[\textcolor{red}{$\blacktriangleright$}] Fine-tuning enables specialization and alignment of LLMs.
    \item[\textcolor{red}{$\blacktriangleright$}] Supervised Fine-Tuning (SFT) is the foundational step.
    \item[\textcolor{red}{$\blacktriangleright$}] RLHF incorporates human preference alignment using PPO.
    \item[\textcolor{red}{$\blacktriangleright$}] DPO streamlines the process with a direct preference loss.
    \item[\textcolor{red}{$\blacktriangleright$}] ORPO and KTO go further: one stage and no reference model, or unpaired thumbs up and down in place of ranked pairs.
    \item[\textcolor{red}{$\blacktriangleright$}] Constitutional AI and RLAIF move the labeling onto a model, trading annotation cost for values inherited from that labeler.
    \item[\textcolor{red}{$\blacktriangleright$}] LoRA and QLoRA make adaptation efficient and accessible; soft prompts, (IA)$^3$ and DoRA trade accuracy against serving cost differently.
    \item[\textcolor{red}{$\blacktriangleright$}] Task vectors combine finished fine-tunes by arithmetic, and are the cheapest answer to a model that has to keep changing.
    \item[\textcolor{red}{$\blacktriangleright$}] Which method to reach for is largely a budget question: SFT to add a capability, preference tuning to shape behaviour, LoRA or QLoRA when memory is the binding constraint.
  \end{itemize}
\end{frame}
